Agent benchmarks for repository-grounded tasks are typically reduced to a
single global success rate, which can obscure where models actually succeed or
fail. We argue that richer structure is already available in existing benchmark
outputs without requiring new metrics or judges. Using \gittaskbench as a case
study, we audit its official \processmetric/\resultmetric split---where
\processmetric marks whether a trace reaches grading and \resultmetric marks
whether it satisfies the task criterion---across a comparable open-model cohort,
supplementary provider-model screens, and full-benchmark provider sweeps. Four
empirical patterns emerge. First, official success is narrow and domain-skewed:
observed \resultmetric mass concentrates in office-document, security, web, and
speech families, while image and video families produce no observed success.
Second, models that tie on global \resultmetric totals succeed on materially
different task families, revealing hidden specialization that a single score
cannot surface. Third, provider models do not form one uniform control
condition: on the same benchmark tasks, \texttt{gpt-5-chat},
\texttt{gpt-5.2}, and \texttt{Gemini-2.5-Pro} activate different families, and
targeted screens can both overstate and miss activations that appear only in
full-benchmark sweeps. Fourth, the intermediate
\processmetric$=$True, \resultmetric$=$False regime is not random residue; it
clusters in structured-extraction families and mixes threshold misses, parser
failures, and evaluation-side errors. These findings support a practical
reporting recommendation: repo-grounded agent evaluations should report the
official \processmetric/\resultmetric split and family-level slices alongside
any global pass rate, so that future method improvements can be attributed to
specific capability gains rather than opaque score movements.
